\section{\textbf{Knowledge Discovery as a Pathway to Knowledge Invention}}
\label{sec:knowledge_discovery}

Knowledge fusion derives a consequential relation by connecting bodies of
knowledge that were previously separated. Knowledge discovery begins from a
different epistemic condition: some feature of reality is not adequately
represented in existing knowledge, even after known information has been
assembled. The missing object may be a phenomenon, entity, mechanism,
variable, regularity, law, material, algorithm, or causal relation. A discovery
system must therefore do more than reorganize descriptions. It must establish
that reality contains a structure that collective human knowledge did not
previously capture.

Machine-assisted discovery is already an active scientific program. Symbolic
regression has recovered compact physical relations from experimental data
\cite{schmidt2009laws,udrescu2020aifeynman}; sparse identification has extracted
governing dynamical equations from measurements \cite{brunton2016sindy}; and
representation-learning systems have proposed latent state variables from raw
observations \cite{chen2022variables}. In biology and chemistry, machine
learning has contributed to protein-structure prediction
\cite{jumper2021alphafold}, antibiotic identification
\cite{stokes2020antibiotic}, materials prediction
\cite{merchant2023materials}, and closed-loop experimental synthesis
\cite{szymanski2023autonomous}. These achievements establish that machines can
participate in multiple stages of discovery. The present framework asks when
that participation becomes machine-originated knowledge invention.

\subsection{\textbf{Prediction, Detection, and Discovery}}

Prediction is not identical to discovery. A system may accurately map inputs
to outputs without identifying a new entity, mechanism, or explanatory
structure. Detection is also insufficient. A detector may locate an anomaly or
classify an unfamiliar object while leaving its significance unexplained.
Optimization can improve an established design without changing the relevant
knowledge. These operations become components of discovery only when they
support a justified claim about something previously unknown.

Let $o$ denote an observation, $p$ a prediction, $a$ a detected anomaly, and
$\kappa_D$ a candidate discovery. The non-equivalence is

\begin{equation}
\begin{aligned}
\operatorname{Predict}_{x}(p\mid o)
&\not\Rightarrow \operatorname{Discover}_{x}(\kappa_D),\\
\operatorname{Detect}_{x}(a\mid o)
&\not\Rightarrow \operatorname{Discover}_{x}(\kappa_D),\\
\operatorname{Optimize}_{x}(y)
&\not\Rightarrow \operatorname{Discover}_{x}(\kappa_D).
\end{aligned}
\label{eq:prediction_not_discovery}
\end{equation}

A prediction may contribute evidence; an anomaly may motivate a hypothesis;
and optimization may expose an unexpected solution. Discovery nevertheless
requires an epistemic transition from an output to a validated claim about the
world or about a formally defined possibility space.

This distinction is illustrated by data-driven law discovery. Sparse
identification methods can produce parsimonious governing equations from noisy
measurements \cite{brunton2016sindy}, while symbolic-regression systems can
recover interpretable relations rather than only point predictions
\cite{schmidt2009laws,udrescu2020aifeynman}. Their explanatory representations
move closer to discovery because they state a general relation that can be
tested outside the observations used to construct it. Even then, recovery of a
known law is rediscovery or methodological validation, not an extension of the
human knowledge frontier.

\subsection{\textbf{Objects of Machine Discovery}}

Machine discovery may target several types of epistemic object. A
\emph{phenomenon} is an observed behavior or effect not adequately represented
by existing theory. An \emph{entity} is a previously unknown object, class,
material, biological structure, or computational construction. A
\emph{mechanism} explains how interacting processes produce an effect. A
\emph{variable} supplies a representation required to describe a system. A
\emph{regularity} expresses a stable relation among observations, and a
\emph{law} provides a generalizable mathematical or causal structure. An
\emph{algorithm} is a previously unknown procedure that produces a formally or
empirically verifiable result.

These objects require different validation standards. A candidate material
must be computationally and ultimately experimentally assessed; a natural law
must predict across conditions; a causal mechanism must survive intervention;
and an algorithm may be established through proof, exhaustive verification, or
implementation benchmarks. Machine-generated algorithmic results demonstrate
the importance of this distinction. AlphaTensor discovered provably correct
matrix-multiplication algorithms that improved on known procedures
\cite{fawzi2022alphatensor}, and AlphaDev produced new sorting routines that
were validated and incorporated into a widely used standard library
\cite{mankowitz2023alphadev}. These are genuine algorithmic discoveries within
human-formulated problem spaces. Under the present account, they demonstrate a
major component of knowledge invention but not yet the autonomous exposure of
an unknown unknown.

Let $\Omega$ denote the relevant space of possible objects and
$\mathcal{R}_{H,t}(\Omega)$ the region represented in collective human
knowledge at time $t$. A candidate machine discovery occupies

\begin{equation}
\begin{aligned}
\kappa_D &\in \Omega,\\
\kappa_D &\notin \mathcal{R}_{H,t}(\Omega),\\
\operatorname{Represent}_{x}(\kappa_D)
&=1.
\end{aligned}
\label{eq:discovery_object}
\end{equation}

The representation condition requires the machine to distinguish the object
from noise and express it in a form that supports examination. It does not yet
establish that the representation is true, important, or attributable to the
machine.

\subsection{\textbf{The Discovery Cycle}}

Knowledge discovery requires a connected cycle rather than a single model
output. The cycle begins with observations and proceeds through anomaly or
structure detection, hypothesis formation, explanatory modeling, discriminating
tests, and validation:

\begin{equation}
\begin{aligned}
\mathcal{O}_{t}
&\longrightarrow \mathcal{A}_{x}\\
&\longrightarrow \mathcal{H}_{x}\\
&\longrightarrow \mathcal{M}_{x}\\
&\longrightarrow \mathcal{T}_{x}\\
&\longrightarrow \mathcal{V}(\kappa_D),
\end{aligned}
\label{eq:discovery_cycle}
\end{equation}

where $\mathcal{O}_{t}$ contains observations, $\mathcal{A}_{x}$ contains
machine-identified anomalies or structures, $\mathcal{H}_{x}$ contains
hypotheses, $\mathcal{M}_{x}$ contains explanatory models,
$\mathcal{T}_{x}$ contains discriminating tests, and
$\mathcal{V}(\kappa_D)$ is the validation result.

The cycle need not be linear. Failed tests may revise the representation,
generate alternative hypotheses, expose measurement error, or cause the system
to seek new observations. Autonomous laboratories demonstrate how prediction,
planning, experimentation, interpretation, and revision can be connected in a
closed loop \cite{boiko2023autonomous,szymanski2023autonomous}. Earlier robot
scientists likewise showed that machine-generated hypotheses could be selected
and tested experimentally \cite{king2004robot,king2009automation}. These systems
provide architectural precedents for the discovery cycle, although their
research domains and objectives were established by humans.

For discovery beyond the human frontier, the machine must exercise epistemic
control across the cycle. It must determine which observation is anomalous,
which hypothesis is explanatory rather than merely predictive, which test can
distinguish competing accounts, and how the result changes the candidate
knowledge claim. A pipeline whose scientific choices are supplied entirely by
humans may automate discovery work without realizing machine discovery in this
strong sense.

\subsection{\textbf{From Anomaly to Explanation}}

Anomalies are important because they may signal that an existing model is
incomplete, but most anomalies are not discoveries. They may result from noise,
instrument error, data leakage, distribution shift, implementation defects, or
unrepresented confounders. A discovery system must discriminate among these
possibilities before attributing an anomaly to a new feature of reality.

Let $a$ be an anomaly under current model $M_H$, and let
$\{h_1,\ldots,h_n\}$ be competing explanations. The system should select a test
$\tau^{*}$ according to its expected ability to discriminate among them:

\begin{equation}
\begin{aligned}
\tau^{*}
&=\arg\max_{\tau\in\mathcal{T}}
\operatorname{Disc}
(\tau;h_1,\ldots,h_n),\\
\operatorname{Disc}
(\tau;h_1,\ldots,h_n)
&>\delta,
\end{aligned}
\label{eq:discriminating_test}
\end{equation}

where $\operatorname{Disc}$ measures the expected discriminating value and
$\delta$ is a domain-appropriate threshold. Strategic experiment selection has
been proposed as a means of accelerating collective scientific discovery
\cite{rzhetsky2015experiments}. Within the present framework, the stronger
requirement is that the machine select tests in service of a hypothesis it
causally originated about a previously unrecognized object.

Explanation also protects against mistaking correlation for discovery. A model
may predict an outcome from a proxy variable while failing under intervention
or distributional change. A candidate discovery should therefore identify
conditions under which its explanation would fail. Falsifiable consequences,
counterfactual behavior, and out-of-distribution predictions provide stronger
evidence than retrospective fit alone.

\subsection{\textbf{Discovery-Driven Knowledge Invention}}

A validated discovery becomes knowledge invention when the machine originates
both the discovery object and the epistemic structure through which it enters
collective knowledge. Let $\operatorname{DKI}_{x}(\kappa_D)$ denote
discovery-driven knowledge invention. Then

\begin{equation}
\begin{aligned}
\operatorname{DKI}_{x}(\kappa_D)
\Longleftrightarrow{}&
\kappa_D\notin\mathcal{K}_{H,t}\\
&\land q_D\notin\mathcal{Q}_{H,t}\\
&\land \operatorname{Discover}_{x}(\kappa_D)\\
&\land \operatorname{Originate}_{x}(\kappa_D\mid W)\\
&\land \operatorname{Explain}_{x}
(\kappa_D,\mathcal{K}_{H,t})\\
&\land \operatorname{Validate}(\kappa_D)=1\\
&\land \operatorname{Consequential}(\kappa_D)=1.
\end{aligned}
\label{eq:discovery_knowledge_invention}
\end{equation}

The first two conditions place the object and its motivating question beyond
current human knowledge and inquiry. The discovery condition requires the
machine to identify the object from evidence or formal search. The origination
condition assigns the causally essential epistemic step to the machine. The
explanation condition relates the object to existing knowledge, while
validation and consequentiality prevent false, trivial, or uninterpretable
novelty from being classified as knowledge invention.

This definition separates three levels of achievement:

\begin{equation}
\begin{aligned}
\operatorname{CandidateDiscovery}
&\subset
\operatorname{ValidatedDiscovery}\\
&\subset
\operatorname{DiscoveryDrivenKnowledgeInvention}.
\end{aligned}
\label{eq:discovery_levels}
\end{equation}

A candidate discovery is an unvalidated representation or hypothesis. A
validated discovery establishes that the object or relation is real within the
relevant standard. Discovery-driven knowledge invention additionally requires
machine origination beyond a human-recognized research target and integration
of the result into an explanatory epistemic structure.

\subsection{\textbf{Failure Modes and Validation Requirements}}

Machine discovery is vulnerable to systematic errors. \emph{Data artifacts}
can masquerade as phenomena. \emph{Overfitting} can produce laws that describe
observations but fail elsewhere. \emph{Proxy discovery} can identify a
predictive correlate instead of a mechanism. \emph{Search-space confinement}
can make an output appear unprecedented even though the system explored only a
human-specified formal space. \emph{Measurement dependence} can create entities
that disappear under a different instrument or preprocessing method.
\emph{Rediscovery} can reproduce knowledge absent from the system's immediate
context but present elsewhere in the literature. Finally,
\emph{validation capture} occurs when the originating model also defines and
evaluates the only test of its claim.

For candidate discovery $\kappa_D$, the evidentiary package should contain

\begin{equation}
\begin{aligned}
\mathcal{E}(\kappa_D)=\{&
\operatorname{ObservationProvenance},
\operatorname{NoveltySearch},\\
&\operatorname{AlternativeHypotheses},
\operatorname{DiscriminatingTests},\\
&\operatorname{Replication},
\operatorname{BoundaryConditions},
\operatorname{Validation}\}.
\end{aligned}
\label{eq:discovery_evidence}
\end{equation}

Observation provenance preserves the origin and transformations of the data.
Novelty search tests whether the discovery already exists in collective
knowledge. Alternative hypotheses prevent premature commitment.
Discriminating tests separate the proposed account from competitors.
Replication assesses whether the result survives independent implementation,
measurement, or experimentation. Boundary conditions specify where the claim
should and should not hold. Validation records the evidence and standards used
to accept the contribution.

Knowledge discovery therefore identifies what reality or a formal possibility
space contains beyond existing representation. It becomes a pathway to
superintelligence only when a generally intelligent machine originates an
unknown object of inquiry, constructs and tests an explanation, and produces a
validated, consequential addition to collective human knowledge. Knowledge
fusion and knowledge discovery can also interact: fusion may reveal the
question that discovery answers, while discovery may supply the missing object
that makes a deeper fusion possible. The next section integrates these pathways
under the paper's central concept of knowledge invention.
